
\documentclass[11pt]{article}
\usepackage[utf8]{inputenc}
\usepackage[T1]{fontenc}
\usepackage[a4paper,margin=1.8cm]{geometry}
\usepackage{array}
\usepackage{booktabs}
\usepackage{longtable}
\usepackage{tabularx}

\setlength{\parindent}{0pt}
\setlength{\parskip}{6pt}
\renewcommand{\arraystretch}{1.18}

\begin{document}

\begin{titlepage}
\centering
{\Large Universidad de las Fuerzas Armadas\par}
{\LARGE \textbf{ESPE}\par}
\vspace{0.35cm}
{\large Departamento de Ciencias de la Computación\par}
\vspace{1.5cm}
{\Large \textbf{Aseguramiento de la Calidad de Software}\par}
\vspace{0.5cm}
{\large \textbf{Grupo: 6}\par}
\vspace{1.2cm}
{\Large \textbf{Informe Completo de Entregables ISO/IEC 25022}\par}
\vfill
\begin{tabular}{rl}
\textbf{Integrantes:} & Caetano Flores, Jordan Guamán,\\
& Anthony Morales, Leonardo Narváez\\
\textbf{NRC:} & 30733\\
\textbf{Docente:} & Ing. Diego Gamboa\\
\end{tabular}
\vfill
{\large 2026\par}
\end{titlepage}


\section*{Introducción}
El presente informe consolida los entregables del taller de medición de Calidad en Uso mediante ISO/IEC 25022 para el caso MediSalud HIS. Se incluyen las fichas de usuario-tarea-contexto, la matriz de priorización, el catálogo de métricas, la obtención de datos, la automatización de indicadores y el plan de mejora continua.

\newpage
\section*{Escenario 4: Identificación de Atributos de Calidad en Uso}

\small
\begin{longtable}{p{0.16\textwidth} p{0.13\textwidth} p{0.23\textwidth} p{0.21\textwidth} p{0.21\textwidth}}
\toprule
\textbf{Proceso} & \textbf{Usuario} & \textbf{Tarea representativa} & \textbf{Contexto de uso} & \textbf{Atributos} \\
\midrule
\endfirsthead
\toprule
\textbf{Proceso} & \textbf{Usuario} & \textbf{Tarea representativa} & \textbf{Contexto de uso} & \textbf{Atributos} \\
\midrule
\endhead
Atención médica y registro de HCE & Médico tratante & Registrar una nota de evolución clínica completa con diagnóstico, indicaciones y firma electrónica. & Consulta externa, 10:00-12:00, red interna de Quito, alta concurrencia. & Eficiencia: tiempo de tarea. Efectividad: completitud. Satisfacción: fluidez. Libertad de Riesgo: información clínica correcta. \\
Agendamiento de citas & Paciente & Agendar una cita con especialista seleccionando sede, fecha, horario y confirmación. & Portal web o app móvil, red celular, 18:00-21:00, sedes de alta demanda. & Efectividad: tasa de éxito. Eficiencia: pasos y tiempo. Satisfacción: confianza. Cobertura: web y móvil. \\
Facturación con seguro médico & Personal de admisión/facturación & Generar factura, validar convenio, procesar pago/copago y emitir comprobante correcto. & Punto de admisión, cierre de jornada, integración con sistema financiero heredado. & Libertad de Riesgo: errores financieros. Efectividad: exactitud. Eficiencia: tiempo de proceso. Cobertura: sede y rol. \\
\bottomrule
\end{longtable}
\normalsize

\subsection*{Conclusión}
La calidad en uso debe medirse sobre tareas concretas, no sobre módulos abstractos. Estas fichas permiten convertir procesos críticos en métricas comparables y accionables.
\newpage
\section*{Escenario 5: Mapeo de Características de Calidad}

\small
\begin{longtable}{p{0.25\textwidth} p{0.10\textwidth} p{0.12\textwidth} p{0.35\textwidth} p{0.08\textwidth}}
\toprule
\textbf{Tarea} & \textbf{Impacto} & \textbf{Frecuencia} & \textbf{Características ISO/IEC 25022} & \textbf{Prioridad} \\
\midrule
\endfirsthead
\toprule
\textbf{Tarea} & \textbf{Impacto} & \textbf{Frecuencia} & \textbf{Características ISO/IEC 25022} & \textbf{Prioridad} \\
\midrule
\endhead
Registrar nota de evolución clínica & Alto & Alta & Eficiencia, Efectividad, Satisfacción, Libertad de Riesgo & 1 \\
Agendar cita en portal & Alto & Alta & Efectividad, Eficiencia, Satisfacción, Cobertura de Contexto & 1 \\
Facturar consulta con seguro & Alto & Media & Libertad de Riesgo, Efectividad, Eficiencia & 2 \\
Completar teleconsulta & Alto & Media & Efectividad, Cobertura de Contexto, Satisfacción & 2 \\
Dispensar medicamento desde receta electrónica & Alto & Media & Libertad de Riesgo, Efectividad & 2 \\
Consultar historial de laboratorio & Medio & Alta & Efectividad, Satisfacción & 3 \\
Generar reporte gerencial mensual & Medio & Baja & Eficiencia, Efectividad & 3 \\
Actualizar datos personales desde app móvil & Bajo & Media & Cobertura de Contexto, Satisfacción & 4 \\
\bottomrule
\end{longtable}
\normalsize

\subsection*{Alcance inicial}
Se recomiendan para medición inmediata las prioridades 1 y 2: HCE, portal de citas, facturación, teleconsulta y dispensación de medicamentos.

\subsection*{Conclusión}
La priorización evita medir todo sin foco y concentra recursos en los procesos con mayor riesgo clínico, financiero y reputacional.
\newpage
\section*{Escenario 6: Catálogo de Métricas ISO/IEC 25022}

\small
\begin{longtable}{p{0.23\textwidth} p{0.16\textwidth} p{0.24\textwidth} p{0.11\textwidth} p{0.10\textwidth} p{0.12\textwidth}}
\toprule
\textbf{Métrica} & \textbf{Característica} & \textbf{Fórmula} & \textbf{Unidad} & \textbf{Umbral} & \textbf{Fuente} \\
\midrule
\endfirsthead
\toprule
\textbf{Métrica} & \textbf{Característica} & \textbf{Fórmula} & \textbf{Unidad} & \textbf{Umbral} & \textbf{Fuente} \\
\midrule
\endhead
Completitud de registro de HCE & Efectividad & Notas completadas / notas intentadas & Proporción & >= 0.95 & logs\_hce.csv \\
Tiempo promedio de registro de HCE & Eficiencia & Suma de tiempos / número de notas & Segundos & <= 8 s & logs\_hce.csv \\
CSAT normalizado & Satisfacción & Promedio CSAT / 5 & Proporción & >= 0.80 & encuesta\_satisfaccion.csv \\
Tasa de errores de facturación & Libertad de Riesgo & Errores de facturación / transacciones & Proporción & <= 0.01 & incidentes\_2025.csv \\
Consistencia entre sedes & Cobertura de Contexto & Mejor tiempo promedio / peor tiempo promedio & Proporción & >= 0.85 & logs\_hce.csv \\
\bottomrule
\end{longtable}
\normalsize

\subsection*{Interpretación}
Los umbrales se fijan antes del cálculo para evitar interpretaciones subjetivas. Cada métrica tiene fuente, unidad y responsable, por lo que puede automatizarse.
\newpage
\section*{Escenario 7: Obtención y Validación de Datos}

\small
\begin{longtable}{p{0.26\textwidth} p{0.12\textwidth} p{0.28\textwidth} p{0.28\textwidth}}
\toprule
\textbf{Archivo} & \textbf{Registros} & \textbf{Origen} & \textbf{Uso} \\
\midrule
\endfirsthead
\toprule
\textbf{Archivo} & \textbf{Registros} & \textbf{Origen} & \textbf{Uso} \\
\midrule
\endhead
data/logs\_hce.csv & 3150 & Generador de logs HCE & Efectividad, Eficiencia y Cobertura de Contexto \\
data/encuesta\_satisfaccion.csv & 150 & Encuesta CSAT simulada & Satisfacción por sede y rol \\
data/incidentes\_2025.csv & 3000 & Conversión desde clasificación Markdown & Libertad de Riesgo y análisis de incidentes \\
\bottomrule
\end{longtable}
\normalsize


\small
\begin{longtable}{p{0.24\textwidth} p{0.10\textwidth} p{0.13\textwidth} p{0.35\textwidth} p{0.10\textwidth}}
\toprule
\textbf{Archivo} & \textbf{Nulos} & \textbf{Duplicados} & \textbf{Rangos inválidos} & \textbf{Resultado} \\
\midrule
\endfirsthead
\toprule
\textbf{Archivo} & \textbf{Nulos} & \textbf{Duplicados} & \textbf{Rangos inválidos} & \textbf{Resultado} \\
\midrule
\endhead
logs\_hce.csv & 0 & 0 & 0 tiempos fuera de 0-120 s; 0 valores inválidos en completada & Válido \\
encuesta\_satisfaccion.csv & 0 & 0 & 0 puntajes fuera de 1-5 & Válido \\
incidentes\_2025.csv & 0 & 0 & No aplica & Válido \\
\bottomrule
\end{longtable}
\normalsize

\subsection*{Resumen}
El tiempo promedio HCE fue 7.43 s y el CSAT promedio fue 3.61/5. Los datos quedaron listos para automatización.
\newpage
\section*{Escenario 8: Automatización de la Medición}

\small
\begin{longtable}{p{0.22\textwidth} p{0.30\textwidth} p{0.38\textwidth}}
\toprule
\textbf{Componente} & \textbf{Archivo} & \textbf{Función} \\
\midrule
\endfirsthead
\toprule
\textbf{Componente} & \textbf{Archivo} & \textbf{Función} \\
\midrule
\endhead
Conversión de incidentes & convertir\_incidentes\_markdown.py & Reconstruye incidentes\_2025.csv desde el material local. \\
Generación de logs & generar\_logs\_hce.py & Simula eventos de registro clínico por sede. \\
Generación CSAT & generar\_encuesta\_satisfaccion.py & Crea encuesta de satisfacción por sede y rol. \\
Validación & validar\_datos.py & Verifica nulos, duplicados y rangos. \\
Cálculo de métricas & metricas\_iso25022.py & Calcula indicadores ISO/IEC 25022. \\
Exportación & exportar\_reporte.py & Genera indicadores.json para dashboards. \\
\bottomrule
\end{longtable}
\normalsize


\small
\begin{longtable}{p{0.17\textwidth} p{0.38\textwidth} p{0.12\textwidth} p{0.12\textwidth} p{0.12\textwidth}}
\toprule
\textbf{Característica} & \textbf{Métrica} & \textbf{Valor} & \textbf{Umbral} & \textbf{Estado} \\
\midrule
\endfirsthead
\toprule
\textbf{Característica} & \textbf{Métrica} & \textbf{Valor} & \textbf{Umbral} & \textbf{Estado} \\
\midrule
\endhead
Efectividad & Completitud de registro de HCE & 0.9651 & 0.95 & Cumple \\
Eficiencia & Tiempo promedio de registro de HCE & 7.43 & 8.0 & Cumple \\
Eficiencia & Notas HCE registradas en 8 segundos o menos & 0.6787 & 0.9 & No cumple \\
Satisfaccion & Indice de satisfaccion CSAT normalizado & 0.7227 & 0.8 & No cumple \\
Libertad de Riesgo & Tasa de errores de facturacion con riesgo economico & 0.0301 & 0.01 & No cumple \\
Cobertura de Contexto & Consistencia de eficiencia entre sedes & 0.9801 & 0.85 & Cumple \\
\bottomrule
\end{longtable}
\normalsize

\subsection*{Interpretación}
Aunque el promedio de HCE cumple 8 s, el RNF-01 no cumple porque solo 67.87\% de notas se registran dentro del umbral. También no cumplen satisfacción y riesgo de facturación.
\newpage
\section*{Informe Ejecutivo y Plan de Mejora Continua}

\small
\begin{longtable}{p{0.25\textwidth} p{0.30\textwidth} p{0.35\textwidth}}
\toprule
\textbf{Hallazgo} & \textbf{Evidencia} & \textbf{Impacto} \\
\midrule
\endfirsthead
\toprule
\textbf{Hallazgo} & \textbf{Evidencia} & \textbf{Impacto} \\
\midrule
\endhead
Incidentes de efectividad concentrados & 1.493 de 3.000 incidentes & Usuarios no completan tareas o las completan con errores. \\
Riesgo clínico y financiero & 721 incidentes de Libertad de Riesgo & Posibles errores de medicación, exposición de datos y problemas de facturación. \\
RNF-01 incumplido & 67.87\% de notas en <= 8 s frente a umbral de 90\% & Lentitud en consulta externa y menor productividad médica. \\
RNF-03 incumplido & 3.01\% de errores de facturación frente a umbral de 1\% & Riesgo financiero y pérdida de confianza. \\
Satisfacción baja & CSAT normalizado de 0.7227 frente a 0.80 & Experiencia percibida insuficiente. \\
\bottomrule
\end{longtable}
\normalsize


\small
\begin{longtable}{p{0.14\textwidth} p{0.34\textwidth} p{0.22\textwidth} p{0.20\textwidth}}
\toprule
\textbf{Fase PDCA} & \textbf{Acciones} & \textbf{Responsable} & \textbf{Evidencia} \\
\midrule
\endfirsthead
\toprule
\textbf{Fase PDCA} & \textbf{Acciones} & \textbf{Responsable} & \textbf{Evidencia} \\
\midrule
\endhead
Planificar & Definir umbrales, tareas prioritarias y fuentes oficiales. & Calidad, TI, Dirección Médica & Catálogo de métricas. \\
Hacer & Ejecutar pipeline semanal y generar indicadores. & TI y Calidad & indicadores.json. \\
Verificar & Comparar resultados contra RNF-01, RNF-03 y CSAT. & Calidad y Auditoría & Informe mensual. \\
Actuar & Optimizar HCE, facturación, UX e integraciones críticas. & TI, Producto, Operaciones & Backlog de mejoras. \\
\bottomrule
\end{longtable}
\normalsize

\subsection*{Recomendación}
Aprobar un ciclo de mejora de 90 días centrado en HCE y facturación, con seguimiento quincenal de indicadores y reporte ejecutivo mensual.


\section*{Conclusión General}
El programa propuesto permite pasar de percepciones aisladas sobre el funcionamiento de MediSalud HIS a evidencia objetiva, repetible y accionable. Los resultados muestran que la organización debe priorizar tres frentes: cumplimiento real del RNF-01 en HCE, reducción de errores de facturación asociados al RNF-03 y mejora de la satisfacción de usuarios clínicos y pacientes.


\end{document}
